To better understand the remaining errors, we examined the first 20 misclassified test examples saved for each of the six runs. Rather than looking at isolated mistakes, we focused on examples that reappeared across multiple models and sequence lengths. This makes it possible to identify systematic weaknesses instead of noise from a single run. Full listings of all misclassified examples are provided in Appendix~\ref{app:misclass}.

The clearest pattern is confusion between \textit{Sci/Tech} and \textit{Business}. Across all six runs, \textit{Sci/Tech} articles misclassified as \textit{Business} is the single most common error type. This is unsurprising: many technology articles in AG News employ business language such as \textit{market}, \textit{shares}, \textit{workers}, or company names. Examples that recur persistently across nearly every run include ``IBM to hire even more new workers'' (App.~\ref{app:run1}, Ex.~6; App.~\ref{app:run3}, Ex.~5; App.~\ref{app:run4}, Ex.~4), ``Bill Clinton Helps Launch Search Engine'' (App.~\ref{app:run1}, Ex.~20; App.~\ref{app:run2}, Ex.~20; App.~\ref{app:run3}, Ex.~17; App.~\ref{app:run5}, Ex.~19), and ``Intuit Posts Wider Loss After Charge'' (App.~\ref{app:run1}, Ex.~17; App.~\ref{app:run2}, Ex.~16; App.~\ref{app:run3}, Ex.~14; App.~\ref{app:run5}, Ex.~15; App.~\ref{app:run6}, Ex.~17). The reverse confusion also occurs: Business articles mislabeled as Sci/Tech include ``Intel to delay product aimed for high-definition TVs'' (App.~\ref{app:run1}, Ex.~10; App.~\ref{app:run2}, Ex.~9; App.~\ref{app:run3}, Ex.~9; App.~\ref{app:run4}, Ex.~9; App.~\ref{app:run5}, Ex.~6; App.~\ref{app:run6}, Ex.~7) and ``Yahoo! Ups Ante for Small Businesses'' (App.~\ref{app:run1}, Ex.~11; App.~\ref{app:run2}, Ex.~10; App.~\ref{app:run3}, Ex.~10; App.~\ref{app:run4}, Ex.~11; App.~\ref{app:run5}, Ex.~7; App.~\ref{app:run6}, Ex.~8), both of which are misclassified in every single run. These systematic errors reveal that the boundary between technological innovation and economic reporting is inherently ambiguous in this dataset.

Another recurring pattern is confusion between \textit{World} and \textit{Business}. Articles about politics or international affairs are repeatedly misclassified as \textit{Business} when they mention oil prices, trade, or financial framing. ``Oil prices bubble to record high'' appears in all six runs (App.~\ref{app:run1}, Ex.~12; App.~\ref{app:run2}, Ex.~11; App.~\ref{app:run3}, Ex.~11; App.~\ref{app:run4}, Ex.~12; App.~\ref{app:run5}, Ex.~8; App.~\ref{app:run6}, Ex.~9), as does ``Venezuela Prepares for Chavez Recall Vote'' (App.~\ref{app:run1}, Ex.~8; App.~\ref{app:run2}, Ex.~6; App.~\ref{app:run3}, Ex.~6; App.~\ref{app:run6}, Ex.~5). The article ``Future Doctors, Crossing Borders'' (App.~\ref{app:run1}, Ex.~13; App.~\ref{app:run2}, Ex.~12; App.~\ref{app:run3}, Ex.~13; App.~\ref{app:run4}, Ex.~16; App.~\ref{app:run5}, Ex.~12; App.~\ref{app:run6}, Ex.~12) is also misclassified in every run, which is a particularly interesting case since its content sits between health and education without any obvious economic keywords. These mistakes suggest that both models rely heavily on common economic keywords to identify the Business class, which leads to confusion when those keywords appear in political or international news contexts.

We also observe a consistent confusion between \textit{World} and \textit{Sports}, particularly around Olympics-related articles. ``Live: Olympics day four'' is misclassified as \textit{Sports} in runs 1, 2, and 3 (App.~\ref{app:run1}, Ex.~9; App.~\ref{app:run2}, Ex.~8; App.~\ref{app:run3}, Ex.~8), despite being labelled \textit{World}. Similarly, ``Afghan women make brief Olympic debut'' (App.~\ref{app:run1}, Ex.~14; App.~\ref{app:run2}, Ex.~13) and ``Security scare as intruder dives in'' (App.~\ref{app:run4}, Ex.~15; App.~\ref{app:run5}, Ex.~10; App.~\ref{app:run6}, Ex.~10) are repeatedly predicted as the wrong class. These cases suggest that Olympic coverage, which blends international news framing with sporting events, is difficult for both architectures. Conversely, some Sports articles such as ``Hamilton Wins Cycling Time Trial Event'' (App.~\ref{app:run2}, Ex.~17; App.~\ref{app:run4}, Ex.~19) are misclassified as \textit{World}, likely because the tour de France context and the athlete's name carry little explicit sporting vocabulary.

Some errors are caused by lexical ambiguity in short texts. ``E-mail scam targets police chief'' (App.~\ref{app:run1}, Ex.~2; App.~\ref{app:run2}, Ex.~2; App.~\ref{app:run3}, Ex.~2; App.~\ref{app:run6}, Ex.~2) is repeatedly misclassified as \textit{World}, most likely because words such as \textit{police} and \textit{chief} are strongly associated with political or crime-related world news. Similarly, ``Super ant colony hits Australia'' (App.~\ref{app:run1}, Ex.~4; App.~\ref{app:run2}, Ex.~3; App.~\ref{app:run3}, Ex.~4) is consistently predicted as \textit{World}, despite belonging to \textit{Sci/Tech}: the country name and event-style headline likely dominate over any scientific signal.

Comparing the two models, the overall error types are strikingly similar. The LSTM shares virtually all of the same repeat errors as the CNN, especially the \textit{Sci/Tech}--\textit{Business} confusions, which suggests that additional sequential modelling capacity is not sufficient to resolve these cases. This is consistent with the Results section: even though the CNN performs best overall at seq64 and the LSTM improves at seq256, both models still fail on many of the same boundary cases. Notably, ``Selling Houston Warts and All, Especially Warts'' (App.~\ref{app:run1}, Ex.~19; App.~\ref{app:run2}, Ex.~19; App.~\ref{app:run3}, Ex.~16; App.~\ref{app:run5}, Ex.~18; App.~\ref{app:run6}, Ex.~20) is misclassified in almost every run, and interestingly by different classes depending on the model and sequence length (\textit{World} for CNN, \textit{Sci/Tech} for LSTM), which shows that some examples are just ambiguous, rather than any modelling weakness.

Compared with Assignment 1 (TF-IDF with Linear SVM and Logistic Regression), the overall error profile is very similar. In both assignments, the dominant confusion is still \textit{Sci/Tech} vs.\ \textit{Business}, with \textit{World} vs.\ \textit{Business} as another recurring boundary. This consistency across the modeling techniques suggests that the main source of error is label ambiguity in the dataset. Many headlines combine technological and financial language, or mix political and economic framing. Moving from sparse TF-IDF features to neural sequence models improves scores, but does not remove these fundamental ambiguities.

Overall, the error analysis suggests that the main challenge is separating articles near the boundary between related topics. The distinction between \textit{Sci/Tech} and \textit{Business} remains the hardest for both models, consistent with the class-wise results in Section~\ref{sec:results} where these two categories had the lowest per-class F1 scores.
